\documentclass[12pt,a4paper]{article}

% --- Packages ---
\usepackage[margin=1in]{geometry}
\usepackage{fontspec}          % XeLaTeX/LuaLaTeX font support
\usepackage{parskip}           % paragraph spacing instead of indentation
\usepackage{microtype}         % typographic refinements
\usepackage{hyperref}          % clickable links
\usepackage{xcolor}            % for hyperref colors

% --- Font ---
\setmainfont{TeX Gyre Termes}  % ships with texlive, Times-like

% --- Hyperref setup ---
% ??: to have underline
\hypersetup{
  colorlinks=true,
  linkcolor=black,
  urlcolor=blue!70!black,
  pdfauthor={Reza Handzalah},
  pdftitle={Statement of Purpose},
}

% --- Header ---
\title{\vspace{-2em}Statement of Purpose}
\author{Reza Handzalah}
\date{\today}

\begin{document}
\maketitle
\thispagestyle{empty}

% ======================================================================
% 1. OPENING — hook + thesis
% ======================================================================
Presentasi tugas akhir pada konferensi nasional SENIATI di ITN Malang pada tahun 2019 menggenapkan masa pembelajaran S1 saya.
Pada intinya, publikasi yang saya paparkan merupakan laporan mengenai sebuah simulasi sistem telekomunikasi.
Meskipun skripsi saya termasuk "skripsi yang baik\footnote{Karena skripsi yang baik adalah skripsi yang selesai.}"
dan artikel ilmiahnya berhasil masuk prosiding\footnote{Dengan 1 (satu) \href{https://scholar.google.com/scholar?cites=10594950162121343959}{sitasi} pada saat saya menulis \emph{statement of purpose} ini.} setelah serangkaian proses, sebenarnya pada saat itu saya masih bertanya-tanya.
Bagaimana bisa simulasi dibuat laporannya, di-\emph{review} oleh editor, lalu dicukupkan dengan penyebarluasan melalui konferensi.
\emph{Bukankah yang namanya "simulasi" belum bermanfaat sampai benar-benar divalidasi pada dunia nyata?}

% ======================================================================
% 2. UNDERGRADUATE RESEARCH
% ======================================================================
\vspace{\baselineskip}
??Statement of purpose ini akan ??circle back ke pertanyaan tersebut. Sebelumnya, izinkan saya secara ringkas menjelaskan latar belakang
akademik ??Strata 1 saya.
Ada dua kegiatan yang menurut saya menarik selama kuliah di Universitas Telkom: Program Kreativitas Mahasiswa (PKM) dan bimbingan tugas akhir.
Bersama tim yang beranggotakan 3 orang, saya menuangkan ide karsa cipta (sesuai deskripsi kategori pada waktu itu) ke dalam proposal.
Puncak dari kompetisi ini adalah helatan PKM nasional, yang kami gagal meraihnya karena hanya lolos di administrasi kampus dan gagal di
seleksi tingkat kopertis.
Kedua, tugas akhir skripsi dan segala ketentuan format dan administrasinya.
Skripsi adalah pengalaman pertama banyak mahasiswa, termasuk saya, menulis argumentasi ilmiah sedemikian rupa yang mengambil posisi
terhadap literatur yang sudah ada.

Walaupun skripsi dan PKM-nya tidak memenangkan penghargaan apa-apa, keduanya memberi saya gambaran mengenai
standar luaran keilmiahan di Indonesia.
Kedua kegiatan ini memberi modal latar belakang saya dalam metodologi dan proses penulisan ilmiah.

% ======================================================================
% 3. PROFESSIONAL EXPERIENCE
% ======================================================================
\vspace{\baselineskip}
Setelah lulus program sarjana, saya masuk ke dunia \emph{software
engineering}. Saya mencoba pengalaman di beberapa perusahaan dengan
industri produk berbeda: mulai dari \emph{digital identity}, \emph{e-commerce platform}, hingga
\emph{beauty} FMCG.

Pengalaman bekerja di industri menanamkan pragmatisme mencegah
``\emph{bug} yang mengganggu \emph{production}''.
Bekerja selaras dengan \emph{culture} dan kebiasaan perusahaan melatih saya untuk berpikir menjaga harmoni dengan
\emph{stakeholder} yang ada. \emph{Scientific mindset} dari latar belakang kuliah bukan dibenturkan dengan realitas, melainkan
tertuntun arahnya menjadi sebuah ketertarikan ilmiah "\emph{research interest}".

% ======================================================================
% 4. RESEARCH INTEREST
% ======================================================================
\vspace{\baselineskip}
Riset yang akan saya ajukan menjadi tesis magister sejalan dengan kegiatan ??Badan Riset dan Inovasi Nasional (BRIN)
khususnya ??Pusat Studi Sains Data dan Informasi yang meneliti \emph{user journey} layanan digital oleh pemerintah
Indonesia. 

Rencananya, saya akan menjadi anggota proyek penelitian
"\emph{Digital Service Journey Simulator} untuk Meningkatkan Inklusivitas Layanan Publik Digital"
di bawah Dr. Ekawati Marlina, peneliti BRIN di bidang Pemerintahan Digital.
Dr. Ekawati Marlina bersama dengan beberapa akademisi dari Universitas Indonesia pernah mengeluarkan publikasi di bidang
sains data (??footnote dari scholar: some papers Ekawati Marlina and Universitas Indonesia faculties have co-authored together) 

Andil kontribusi saya terhadap proyek tersebut saya tuliskan dalam bentuk proposal yang berjudul
\emph{"Beyond WCAG Grading: Enumerating Accessibility Metrics in Web Form Interaction Process of
Indonesia E-Government Websites"}.
Proposal ini terlampir dalam pendaftaran resmi saya pada program \href{}{\emph{Degree by Research} oleh BRIN}

Saya juga secara independen mengembangkan \emph{HTTP testing tool} yang terinspirasi saat bekerja di industri.
Penggunaan utama \emph{tool} ini adalah sebagai \emph{load testing tool} yang 
memungkinkan simulasi sistem digunakan banyak orang secara bersamaan.
Walaupun belum tentu akan menjadi \emph{reusable building block}, pengembangan ini saya yakini memiliki tema yang berdekatan
dengan konsep \emph{virtual user} pada penelitian \emph{Digital Service Journey Simulator}.


% ======================================================================
% 5. FIT & CLOSING
% ======================================================================
\vspace{\baselineskip}

Saya tidak tahu apakah pertanyaan mengenai kebermanfaatan simulasi yang saya tulis di paragraf pembuka memiliki jawaban yang definitif dan mutlak benar.
Saya hanya ingin menerapkan metodologi dasar \emph{simulation science}.
Saya percaya simulasi yang \emph{academically rigorous} -lah yang bermanfaat bagi para \emph{engineer}, \emph{architect},
dan penentu kebijakan dalam menjawab tantangan skalabilitas sistem digital apapun yang menopang kebutuhan manusia di era sekarang.

Saya meyakini, \emph{curiosity} di bidang apapun akan berdampak baik bukan hanya terhadap pemilik keingintahuan, tetapi juga
terhadap kekayaan khazanah di bidang tersebut.
Untuk memberikan peluang yang lebih baik bagi keingintahuan saya berbuah sesuatu yang berarti,
saya ingin menjadi bagian dari komunitas bernama institusi pendidikan pascasarjana Teknologi Informasi.
Lebih khusus, saya ingin berkolaborasi dengan lab \href{https://scholar.ui.ac.id/en/organisations/e-government-e-business-egb/}{E-Government \& E-Business (EGB)} Universitas Indonesia.

Terakhir, semoga setiap niat baik kita dimudahkan oleh Yang Maha Kuasa agar memberikan hasil dan manfaat yang baik pula.

\end{document}